\paragraph*{Problems}

\begin{enumerate}[label=\thesection.\arabic*.,ref=\thesection.\arabic*,start=24,series=book-problems]

\item\label{ex:bounded-operator-space-completeness} 证明 Theorem~\ref{thm:bounded-operator-space-completeness}.

\item\label{ex:hahn-banach-hilbert-extension-details} 验证 Section~\ref{subsec:hahn-banach-theorem} 末尾的论述.

\item\label{ex:nondense-linear-manifold-annihilating-functional} 设 $M$ 为 Banach 空间 $X$ 中的线性流形. 证明 $M$ 在 $X$ 中不稠密的充分必要条件是存在非零 $l\in X^*$, 使\[l(x)=0\quad\text{for every }x\in M.\]

\item\label{ex:infinite-hilbert-operator-space-nonseparable} 设 $H$ 为无穷维 Hilbert 空间. 证明 $\mathcal{L}(H)$ 不可分.

\item\label{ex:orthonormal-basis-weak-zero} 验证 Example~\ref{ex:orthonormal-basis-weak-zero-example} 的断言.

\item\label{ex:bounded-linear-map-weak-continuity} 设 $X,Y$ 为 Banach 空间, $T\in\mathcal{L}(X,Y)$, $x_n$ 为 $X$ 中的弱收敛序列. 证明 $Tx_n$ 在 $Y$ 中弱收敛.

\item\label{ex:closed-linear-subspace-weak-sequential-closure} 设 $M$ 为 Banach 空间 $X$ 的闭子空间, $x_n\in M$ 弱收敛到 $x$. 证明 $x\in M$. Hint: 先证明, 若对每个满足 $f|_M=0$ 的 $f\in X^*$ 都有 $f(x)=0$, 则 $x\in M$.

\item\label{ex:hilbert-weak-convergence-and-norms-imply-strong} 设 $H$ 为 Hilbert 空间. 设 $x_n$ 弱收敛到 $x$, 并另假设 $\|x_n\|\to\|x\|$. 证明 $x_n\to x$ 在范数意义下 (强) 成立.

\item\label{ex:bounded-l-two-distribution-limit-is-weak-limit} 设 $f_n$ 为 $L^2(\Omega)$ 中的有界序列, 并假设 $f_n$ 在分布意义下收敛. 证明 $f_n$ 弱收敛.

\item\label{ex:weak-strong-holder-dual-product-distribution-limit} 设 $p\in(1,\infty)$, $q=p/(p-1)$. 假设\[f_n\rightharpoonup\bar f\quad\text{(weakly) in }L^p(\Omega),\]\[g_n\to\bar g\quad\text{(strongly) in }L^q(\Omega).\]证明\[f_ng_n\to\bar f\bar g\quad\text{in }\mathcal{D}^{\prime}(\Omega).\]

\item\label{ex:dirac-not-l-one-and-l-one-nonreflexivity} (a) 证明 $\delta(x)\notin L^1(\mathbb{R})$. Hint: 对给定的 $g\in\widetilde L^1(\mathbb{R})$, 找到连续函数 $f$, 使
\begin{equation}\label{eq:dirac-not-l-one-continuous-test-separation}
\left|\int_{-\infty}^{\infty}f(x)g(x)\,\mathrm{d}x-f(0)\right|\geq\dfrac{1}{2}\max_{x\in\mathbb{R}}|f(x)|.
\end{equation}
(b) 证明 $L^1(\mathbb{R})$ 不自反.

\end{enumerate}
